Este informe presenta un análisis experimental de cuatro algoritmos de ordenamiento (MergeSort, QuickSort, PatienceSort y \texttt{std::sort}) y dos algoritmos de multiplicación de matrices cuadradas (el método iterativo tradicional y el de Strassen), con el objetivo de contrastar su complejidad asintótica con su desempeño medido. Se ejecutaron 288 mediciones de ordenamiento y 144 de multiplicación, cubriendo todos los tamaños, tipos de entrada, dominios y muestras especificados en el enunciado. Los resultados muestran que \texttt{std::sort} obtiene el menor tiempo en 21 de las 24 configuraciones de ordenamiento; que el esquema de partición de Hoare evita la degradación a $\mathcal{O}(N^2)$ que el esquema de Lomuto exhibe ante entradas con pocos valores distintos; y que PatienceSort presenta una sensibilidad extrema al orden inicial, con una diferencia de 51 veces en tiempo y 46 veces en memoria entre su peor caso (entrada ascendente) y su mejor caso (entrada descendente) para $N = 10^7$. En multiplicación de matrices, la ventaja asintótica de Strassen sólo se observa de forma consistente en $N = 1024$, donde resulta entre 1{,}07 y 2{,}01 veces más rápido a costa de 2{,}4 veces más memoria residente.
